% RESEARCH_DRAFT_SHA256: b17861f697fd77ca61e61093b488c497fd2cf1c6f2e3beb75ec4cc96890b656e
% Layout and source-reference wrappers expanded for the web renderer.
\documentclass[11pt]{article}
\usepackage[margin=25mm]{geometry}
\usepackage{amsthm}
\usepackage{fontspec,unicode-math}
\setmainfont{Latin Modern Roman}
\setmathfont{Latin Modern Math}
\usepackage{amsmath,amssymb,mathtools}
\usepackage{xurl,hyperref}
\hypersetup{colorlinks=true,urlcolor=blue,linkcolor=blue}
\setcounter{secnumdepth}{0}
\setlength{\parindent}{0pt}
\setlength{\parskip}{5pt}
\setlength{\emergencystretch}{3em}
\newcommand{\sref}[2]{\hyperlink{source:#1:#2}{[S#2]}}
\newcommand{\cataloguescope}[1]{\paragraph{Source question.}#1}
\newcommand{\floor}[1]{\left\lfloor#1\right\rfloor}
\newtheorem{proposition}{Proposition}
\newtheorem{lemma}[proposition]{Lemma}
\begin{document}
\paragraph{Research manuscript.} OpenAI \texttt{gpt-6-astra}, reasoning \texttt{ultra}; 11 October 2026.
\paragraph{Preserved source body.} The complete catalogue mathematical text and sources follow with layout and source links expanded, including source qualifications and its historical preparation notes. New research begins at the section headed Research attempt.

\section{Three-Terminal Distance-Interdiction Complexity}\label{statement}

\subsection{Definitions and mathematical statement}

\paragraph{Shared notation.}Standard number systems are denoted by $\mathbb N,\mathbb Z,\mathbb Q,\mathbb R,\mathbb C$; index ranges, computational models and specialized notation follow the cited problem source and the definitions below.


$d(u,v)$ is shortest-path length in the original graph; after deletion it is computed in the remaining graph, with disconnected pairs assigned infinite distance. In the unweighted-cost variant, deletion cost is the number of removed edges.
\paragraph{Statement recorded in the supplied source.}
Given an undirected graph, with or without edge costs, and three vertices $r,s,t$, determine the complexity of finding a minimum-cost edge set whose deletion increases each of $d(r,s)$, $d(r,t)$, and $d(s,t)$ by at least $1$. \href{https://doi.org/10.4171/owr/2005/50}{[S2]}
\paragraph{Scope qualification.}
The supplied graph is undirected; the separately mentioned directed variant is not added to its target. Edge lengths and deletion costs should not be conflated.
\paragraph{Source status.}
The supplied catalogue labels this record open. Its source-status metadata was checked against the supplied snapshot on 2026-10-11. This import does not independently establish current open status or mathematical correctness. The catalogue literature assessment is dated 2026-08-21.
\subsection{Short English statement}
Classify minimum-cost edge deletion that increases all three pairwise terminal distances.
\subsection{Sources}
\begin{itemize}
\item[S1] Ulam AI and the UnsolvedMath Contributors, UnsolvedMath: A Curated Collection of Open Mathematics Problems, catalogue entry. Catalogue curation and metadata: CC BY 4.0. The supplied source question is retained; definitions and the English summary are editorial additions.
\url{https://huggingface.co/datasets/ulamai/UnsolvedMath}.
\item[S2] Combinatorial Optimization (2006), Oberwolfach Reports, DOI 10.4171/owr/2005/50, pp. 2833--2932
\url{https://doi.org/10.4171/owr/2005/50}.
\textit{Underlying Oberwolfach Reports material retains its source-specific licence, including CC BY-SA 4.0 where applicable; see THIRD\_PARTY\_NOTICES.md.}
\end{itemize}
\paragraph{Preparation and review.}
Original question attributed to the cited source and UnsolvedMath contributors. Codex (GPT-6; exact serving revision not exposed) assisted with definitions, English summaries, source inspection and formatting. No new human mathematical authorship or independent proof verification is asserted. Catalogue identity checked against the complete supplied registry and existing repository statements; known flagged repetitions and obvious identical conjectures excluded. Full mathematical-equivalence/literature review has not been performed. Statement preparation and final review followed contribution prompts C and D; automated checks verify structure and identity, not open status.
\label{end:statement}

\section{Research attempt}
\subsection{Literature, source context and scope}
The full original catalogue body and source qualifications are preserved above. The relevant primary workshop passages were read in the cached publisher report. Historical claims and subsequent attributed results are distinguished from the mathematical work below. A restricted example, standard lemma or known theorem is not counted as new progress on the entire original target. Missing assumptions and repeated or rhetorical extracts are not silently turned into different solved problems.

\subsection{The three-terminal question on a tree}
The source distinguishes edge lengths from deletion costs and asks about hitting every shortest path for three terminal pairs forming an undirected triangle. We first restrict the underlying graph to a tree, with unit edge lengths and nonnegative deletion costs. For three distinct terminals $a,b,c$, their pairwise paths meet at a unique median vertex $v$; the three terminal-to-$v$ arms are edge-disjoint. Let $m_a,m_b,m_c$ be the minimum edge costs on these arms, with value $+\infty$ for an empty arm.
\begin{proposition}
The minimum cost of increasing all three pairwise distances on the tree is the sum of the two smallest arm minima.
\end{proposition}
\begin{proof}
The union of the three pairwise paths is either a subdivided claw or a path with one terminal between the other two, giving the stated median and arms. Each pairwise path consists of the corresponding two arms. Deleting an edge outside these arms affects none of the three paths. Because paths in a tree are unique, a pair's distance increases exactly when its path is cut; the new distance is then infinite. All three pairwise paths are cut if and only if at least two of the three arms are cut. Necessity follows because if at most one arm is cut, the path joining the other two terminals survives; sufficiency follows because every pair of arms includes one of the two cut arms. Any cut arm costs at least its cheapest edge, and deleting one cheapest edge from each of the two cheapest nonempty arms attains their sum. At most one arm is empty because the terminals are distinct, so two finite minima exist.
\end{proof}
This gives a linear-time tree algorithm after finding the three paths. A three-arm star with costs $1,2,3$ has optimum 3. However the route cannot be extended simply by replacing shortest-path interdiction with terminal separation on arbitrary graphs.
\begin{proposition}
In the unit-length, unit-cost complete graph on terminals $a,b,c$ and a fourth vertex $v$, increasing all three terminal distances costs 3, whereas separating all three terminals into different components costs 5.
\end{proposition}
\begin{proof}
Each terminal pair initially has distance one and its direct edge is its unique length-one path. All three terminal-to-terminal edges must therefore be deleted, and these three deletions suffice: the remaining star through $v$ has all terminal distances two. For complete separation those three edges still must go, and at most one of the three edges from $v$ to a terminal may remain, since two would connect two terminals. Thus at least five edges must be deleted, and keeping exactly one such edge attains five.
\end{proof}
\subsection{Verification and remaining gaps}
All subsets of the star's three edges and of the complete graph's six edges were enumerated. The cost distinction shows the precise failure of transporting the tree multiway-cut description to graphs with alternate longer paths. Neither a polynomial algorithm nor an NP-hardness reduction for the report's general undirected three-pair case is obtained. The source already records easier and harder nearby cases; these toy cases do not change the original complexity status.

\section{References and source attribution}
Primary workshop report: \url{https://doi.org/10.4171/owr/2005/50}. Relevant passages read: 2005/50 pp. 2856--2858, shortest-path interdiction, unit lengths versus costs, and three terminal pairs.

All retrievals were checked on 11 October 2026. Original question attribution remains in the preserved Sources section. External theorems are attributed and not claimed as new discoveries or fully verified merely by citation.


\paragraph{Integration review.} Codex (GPT-6; exact serving revision not exposed) reviewed the written arguments and expanded layout and source-reference wrappers. Repository metadata and the production web renderer were checked using Python 3.13 and JavaScript with MathJax 3.2.2. This model-assisted review is not independent human proof verification; no PDF compilation or formal proof checking is asserted.

\section{Completion Estimate}
\noindent\textbf{COMPLETION: 0\%}
No substantiated new progress toward the entire original target is claimed. The worked results and failed route are delimited above; a known result, illustrative model, or auxiliary obstruction is not counted as original-target completion.

\section{Final status and attribution}
\textbf{Attempt status: unresolved.} This describes this research attempt, not a new ruling on the source catalogue's status or on external literature claims.
The mathematics and this self-audit were generated by OpenAI \texttt{gpt-6-astra}, reasoning \texttt{ultra}, on 11 October 2026. Codex (GPT-6; exact serving revision not exposed) orchestrated the batch and contributed formatting/repository review. No human mathematical authorship, independent human proof checking, or formal verification is claimed. The historical preparation paragraphs in the reproduced catalogue body concern that body only.

Contribution prompt A was applied to the complete definition and supplied source record, followed by the adversarial checks of prompt B. Local reproduction files are \texttt{batch016/checks.py}, \texttt{computations.json}, and the per-ID audit JSON. They were not public links at preparation time. Numerical checks are finite and do not prove unrestricted results. No LaTeX engine was found locally; compilation was not performed by this research agent.

\end{document}
